\section{Decision \& Recommended Model(s)}
\label{sec:decision}

\subsection{Final Decision and Rationale}
\label{subsec:final_decision}

We adopt Google \textbf{Gemini 2.5 Flash} as the primary tutor model for Mettle's Socratic chatbot. This choice follows the multi-criterion weighting in \cref{tab:criteria_weights}, the statistical evidence in \cref{sec:statistical_analysis}, the aggregate outcomes in \cref{sec:results}, and the operational analyses in \cref{sec:experiments,sec:cost_modelling,sec:safety}. Three factors tipped the decision:

\begin{enumerate}
	\item \textbf{Pedagogical superiority with statistical backing.} Gemini 2.5 Flash delivers the highest adjudicated pedagogy score in \cref{tab:summary_scores} (3.91 with a 95\% CI of 3.86--3.96) and achieves decisive Wilcoxon advantages over Command R, Llama 4 Maverick, and GPT-4o mini (Hodges--Lehmann gains of +0.10, +0.26, and +0.045; $q\leq 1.5\times10^{-2}$) in \cref{fig:effect_size_heatmap}. The Friedman statistic in \cref{subsec:comparative_analysis} confirms a real separation ($\chi^2(4)=120.07$, $p=5.2\times10^{-25}$). These results advance the 35\%-weighted pedagogical criterion and align with the strict Socratic persona outcomes validated in \cref{subsubsec:few_shot_results}.
	\item \textbf{Context fidelity without quality trade-offs.} Gemini leads the inclusion leaderboard in \cref{fig:context_inclusion_diff} (94\% inclusion, $q=0.017$ versus Llama), maintains parity with Claude on contextual scores (0.008 HL gain; $q=0.33$), and preserves low hallucination incidence (2\% in \cref{tab:summary_scores}). Its handling of ephemeral context insertion remains stable after the on-the-fly contextualization test in \cref{subsec:contextualization_test}, satisfying the 25\%-weighted contextual adaptability axis.
	\item \textbf{Operational efficiency and safety.} Gemini exhibits the lowest observed latency across the cohort (median 521~ms, 95th percentile 810~ms) with bootstrap intervals excluding zero against all peers in \cref{fig:latency_forest}, directly servicing the latency \& throughput objective. Cost projections in \cref{sec:cost_modelling} keep Gemini within budget for both the classroom and departmental scenarios despite slightly higher per-call pricing than Command R; the token savings from ephemeral context passing (\cref{subsubsec:contextualization_results}) offset the price difference. Safety telemetry mirrors cohort norms: the sanitization proxy in \cref{tab:pii_proxy_counts} reports zero persisted PII and no safety exceptions unique to Gemini, while reliability remains acceptable (93\% first-try in \cref{tab:summary_scores}, reinforced by the ICC and $\kappa$ levels in \cref{subsec:inter_rater_reliability}).
\end{enumerate}

Competing models excel on narrower axes---GPT-4o mini on reliability (98\% first-try) and Command R on inclusion (96\%) and price---yet neither offers Gemini's simultaneous strength across the pedagogical, contextual, and latency dimensions that dominate the scoring rubric. Given the additive decision rule in \cref{eq:overall_score}, Gemini maximizes weighted utility while staying within the deployment cost envelope defined in \cref{subsec:projected_scenarios}.

\subsection{Recommendation for a Hybrid Approach}
\label{subsec:hybrid_approach}

For high-stakes deployments we retain an optional blended routing scheme that complements Gemini 2.5 Flash with Cohere Command R 08-2024. The blend exploits Command R's superior context inclusion (+11.7 percentage points over Llama; \cref{fig:context_inclusion_diff}) and lower marginal cost, while Gemini handles pedagogically demanding or latency-sensitive turns. We parameterize routing with the convex combination in \cref{eq:routing_cost}, selecting $\rho$ (Gemini share) by prompt category. Categories that historically triggered misconception repair or substantial scaffolding---``Socratic Probing'' and ``Feedback \& Improvement'' in \cref{subsec:category_differentiation}---receive Gemini by default, whereas ``Edge Cases'' and routine context confirmations fall back to Command R when latency and cost incentives dominate.

Routing decisions are informed by the harness metadata: the SQLite store curated in \cref{subsubsec:data_persistence} exposes per-category latency and inclusion rates, enabling scripted policies that flag prompts with low prior $Q$ or high hallucination risk for Gemini. Because both adapters share the standardized interface in \cref{subsubsec:adapter_interface} and comply with the sanitization pathway in \cref{subsec:sanitization_procedures}, switching incurs no additional integration complexity. Empirical cost ceilings remain under the departmental target when $\rho\leq0.7$, as shown by substituting the observed token medians into \cref{eq:routing_cost}.

\subsection{Comparative Assessment by Evaluation Axis}
\label{subsec:axis_assessment}

We summarize model standing per evaluation axis from \cref{sec:evaluation_criteria}, grounding each claim in measured artifacts.

\paragraph{Pedagogical quality (35\%).} Gemini leads on $Q$ with statistically supported pairwise wins in \cref{fig:effect_size_heatmap}. Claude's mean is close but does not significantly exceed Gemini ($q=0.32$). GPT-4o mini is competitive but trails on depth-related rubric dimensions noted in \cref{subsec:pedagogical_quality}. Llama and Command R underperform with large negative Cliff's $\delta$ against Gemini.

\paragraph{Contextual adaptability (25\%).} On the continuous fidelity score, Gemini and Claude form the top tier (\cref{subsec:statistical_results}), while Command R maximizes binary inclusion (\cref{fig:context_inclusion_diff}). Ephemeral context passing (\cref{subsec:contextualization_test}) sustained inclusion without degrading hallucination rates; thus, the top tier persisted under deployment-default settings.

\paragraph{Safety \& ethics (10\%).} All models operated within the guardrails of \cref{sec:safety}. The PII proxy (\cref{tab:pii_proxy_counts}) reported zero persistence events; hallucination proxies and human adjudication flagged no model-specific safety regressions.

\paragraph{Cost (10\%).} Command R is the most economical per 1k interactions in \cref{tab:summary_scores}; Gemini's higher quality offsets cost via routing and token savings (\cref{subsec:projected_scenarios,subsec:sensitivity_analysis}).

\paragraph{Latency \& throughput (5\%).} Gemini shows the lowest median and tail latency (\cref{fig:cdf_latency,fig:latency_forest}).

\paragraph{Reliability \& availability (5\%).} GPT-4o mini achieved the highest first-try rate (98\%), while Gemini maintained acceptable reliability (93\%), both within the harness's error-handling envelope (\cref{subsubsec:core_functionality}).

\paragraph{Ease of integration \& tooling (5\%).} All selected providers integrated through a uniform adapter (\cref{subsubsec:adapter_interface}). Gemini's client and instrumentation matched the harness contract without workarounds; Llama required host-specific adjustments.

\paragraph{Controllability \& instruction following (5\%).} The strict Socratic template improved $Q$ across models (\cref{subsubsec:few_shot_results}); temperature $T=0.3$ balanced consistency and depth (\cref{subsubsec:temperature_results}).

\subsection{Sensitivity of the Decision to Weights}
\label{subsec:decision_sensitivity}

Because the overall score (\cref{eq:overall_score}) is a weighted sum, we examined robustness to weight perturbations of \textpm{}10 percentage points on major axes while holding the total at 100\%. Gemini remained the argmax when (i) pedagogy received as little as 25\% and cost as much as 20\%, or (ii) contextual adaptability rose to 35\%. Only an extreme tilt favoring cost ($\geq 30\%$) combined with down-weighted pedagogy ($\leq 25\%$) produced a tie with Command R, at which point the latency gap (\cref{fig:latency_forest}) and mixed qualitative signals broke the tie in Gemini's favor. Thus, the decision is stable under plausible governance-driven reweightings.

\subsection{Limitations of the Decision}
\label{subsec:decision_limitations}

The study used a synthetic, curated prompt bank (\cref{subsec:datasets_prompt_bank}); external validity to live, long-horizon tutoring may differ. Tokenization and pricing assumptions (\cref{sec:cost_modelling}) evolve over time. Finally, human ratings, although reliable (\cref{subsec:inter_rater_reliability}), remain subjective; we therefore anchored the decision in converging quantitative tests and effect sizes (\cref{sec:statistical_analysis}) and retained a hybrid fallback (\cref{subsec:hybrid_approach}) to accommodate drift.
